\chapter{Code Reference}

\section{Project Structure}

\begin{forest}
  for tree={
    font=\ttfamily,
    grow'=0,
    child anchor=west,
    parent anchor=south,
    anchor=west,
    calign=first,
    edge path={
      \noexpand\path [draw, \forestoption{edge}]
      (!u.south west) +(7.5pt,0) |- node[fill,inner sep=1.25pt] {} (.child anchor)\forestoption{edge label};
    },
    before typesetting nodes={
      if n=1
        {insert before={[,phantom]}}
        {}
    },
    fit=band,
    before computing xy={l=15pt},
  }
[Project Root
  [data/
    [raw/
      [Pay\_emploier\_sept25.xlsm]
      [Control\_Panel.xlsm]
    ]
  ]
  [docs/
    [latex/ (LaTeX documentation)]
    [DASHBOARD\_GUIDE.md]
  ]
  [outputs/
    [payroll\_long.csv]
    [forecast\_ar2.csv]
    [results\_*.csv (model results)]
  ]
  [scripts/
    [run\_excel.py]
    [run\_dashboard.py]
    [run\_forecast.py]
    [run\_load\_data.py]
    [run\_plot.py]
  ]
  [src/
    [\_\_init\_\_.py]
    [core/
      [config.py]
      [data\_loader.py]
      [multi\_sheet\_loader.py]
      [normalization.py]
    ]
    [data generation/
      [generate\_synthetic\_payroll.py]
    ]
    [excel\_tools/
      [import\_dialog.py]
      [mapping\_gui.py]
    ]
    [features/
      [analytics/
        [dashboard\_gui.py]
        [plot\_salary\_field.py]
        [regression\_*.py (models)]
        [visualize\_all\_models.py]
      ]
      [gui/
        [excel\_gui\_launcher.py]
      ]
      [payroll\_processing/
        [automation.py]
        [mapping\_engine.py]
      ]
    ]
  ]
  [tests/
    [test\_monthly\_mapping\_engine.py]
    [test\_monthly\_mapping\_gui\_launcher.py]
  ]
  [Control\_Panel\_New.xlsm (Excel interface)]
  [requirements.txt]
]
\end{forest}

\section{Core Modules}

\subsection{src/core/config.py}

\textbf{Purpose:} Centralized configuration for paths, constants, and settings.

\textbf{Key Constants:}

\begin{lstlisting}[style=pythonstyle]
CONFIG = {
    "paths": {
        "project_root": Path(__file__).parent.parent.parent,
        "data_dir": "data/raw",
        "output_dir": "outputs",
    },
    
    "excel": {
        "id_row": 7,              # Row with employee IDs
        "code_col": 0,            # Column A (0-indexed)
        "label_col": 2,           # Column B
        "first_data_col": 2,      # Column C (where employees start)
    },
    
    "output": {
        "payroll_long": "outputs/payroll_long.csv",
        "forecast": "outputs/forecast_ar2.csv",
    },
    
    "month_sheets": [
        "January", "February", "March", "April", 
        "May", "June", "July", "August", 
        "September", "October", "November", "December"
    ],
    
    "custom_sheet": "CustomMap",
    
    "target_columns": [
        "employee_id",
        "salaire_brut",
        "salaire_net",
        "cotisation_salariale",
        "cotisation_patronale",
        # ... more columns ...
    ]
}
\end{lstlisting}

\textbf{Usage:}
\begin{lstlisting}[style=pythonstyle]
from src.core.config import CONFIG

data_path = CONFIG["paths"]["data_dir"]
id_row = CONFIG["excel"]["id_row"]
\end{lstlisting}

\subsection{src/core/normalization.py}

\textbf{Purpose:} Normalize labels for fuzzy matching (remove accents, lowercase, strip whitespace).

\textbf{Key Functions:}

\begin{lstlisting}[style=pythonstyle]
def _norm_label(text: str) -> str:
    """
    Normalize label for matching.
    
    - Lowercase
    - Remove accents
    - Strip whitespace
    - Collapse multiple spaces
    
    Example:
        "Salaire  BASE" -> "salaire base"
        "Cotisation Salariale" -> "cotisation salariale"
    """
    import unicodedata
    
    if not text:
        return ""
    
    # Convert to string
    text = str(text)
    
    # Lowercase
    text = text.lower()
    
    # Remove accents
    text = ''.join(
        c for c in unicodedata.normalize('NFD', text)
        if unicodedata.category(c) != 'Mn'
    )
    
    # Strip and collapse spaces
    text = ' '.join(text.split())
    
    return text
\end{lstlisting}

\textbf{Usage:}
\begin{lstlisting}[style=pythonstyle]
from src.core.normalization import _norm_label

norm1 = _norm_label("SALAIRE BASE")   # "salaire base"
norm2 = _norm_label("Salaire  Base")  # "salaire base"
assert norm1 == norm2  # True - fuzzy match
\end{lstlisting}

\subsection{src/core/data\_loader.py}

\textbf{Purpose:} Load payroll data from various formats (Excel, CSV).

\textbf{Key Functions:}

\begin{lstlisting}[style=pythonstyle]
def load_payroll_long(csv_path: str) -> pd.DataFrame:
    """
    Load long-format payroll CSV.
    
    Returns DataFrame with columns:
        - month (datetime)
        - employee_id (str)
        - salaire_brut (float)
        - salaire_net (float)
        - ...
    """
    df = pd.read_csv(csv_path)
    df['month'] = pd.to_datetime(df['month'])
    return df


def load_mapped_workbook(wb_path: str) -> Dict[str, pd.DataFrame]:
    """
    Load all month sheets from mapped workbook.
    
    Returns:
        Dictionary mapping month name to DataFrame
    """
    import xlwings as xw
    
    wb = xw.Book(wb_path)
    data = {}
    
    for sheet in wb.sheets:
        if sheet.name in CONFIG["month_sheets"]:
            # Read sheet to DataFrame
            used_range = sheet.range("A1").current_region
            df = pd.DataFrame(used_range.value[1:], 
                            columns=used_range.value[0])
            data[sheet.name] = df
    
    wb.close()
    return data
\end{lstlisting}

\subsection{src/core/multi\_sheet\_loader.py}

\textbf{Purpose:} Load and combine multiple month sheets into single long-format DataFrame.

\textbf{Key Functions:}

\begin{lstlisting}[style=pythonstyle]
def load_all_months_to_long(wb_path: str) -> pd.DataFrame:
    """
    Load all month sheets and combine to long format.
    
    Process:
        1. Load each month sheet
        2. Add 'month' column
        3. Concatenate all months
        4. Sort by month and employee_id
    
    Returns:
        Long-format DataFrame
    """
    import xlwings as xw
    from datetime import datetime
    
    wb = xw.Book(wb_path)
    all_data = []
    
    for i, month_name in enumerate(CONFIG["month_sheets"]):
        try:
            sheet = wb.sheets[month_name]
        except:
            continue
        
        # Read sheet
        used_range = sheet.range("A1").current_region
        if used_range.rows.count <= 1:
            continue
        
        df = pd.DataFrame(used_range.value[1:],
                         columns=used_range.value[0])
        
        # Add month column (assuming year 2024)
        df['month'] = datetime(2024, i+1, 1)
        
        all_data.append(df)
    
    wb.close()
    
    # Combine
    combined = pd.concat(all_data, ignore_index=True)
    combined = combined.sort_values(['month', 'employee_id'])
    
    return combined
\end{lstlisting}

\section{Excel Integration Modules}

\subsection{src/excel\_tools/import\_dialog.py}

\textbf{Purpose:} File browser dialog for selecting source workbook.

\textbf{Key Functions:}

\begin{lstlisting}[style=pythonstyle]
def select_workbook_dialog() -> str:
    """
    Open file dialog to select Excel workbook.
    
    Returns:
        Absolute path to selected workbook, or None
    """
    from tkinter import Tk
    from tkinter.filedialog import askopenfilename
    
    root = Tk()
    root.withdraw()  # Hide main window
    
    filepath = askopenfilename(
        title="Select Payroll Workbook",
        filetypes=[
            ("Excel Files", "*.xlsx *.xlsm *.xls"),
            ("All Files", "*.*")
        ]
    )
    
    root.destroy()
    return filepath if filepath else None
\end{lstlisting}

\subsection{src/excel\_tools/mapping\_gui.py}

\textbf{Purpose:} Main GUI for defining and executing mappings.

\textbf{Key Classes:}

\begin{lstlisting}[style=pythonstyle]
class MappingGUI:
    """Main mapping interface."""
    
    def __init__(self, source_wb_path: str):
        self.source_wb_path = source_wb_path
        self.mappings = {}  # {dest_label: formula_string}
        
        # Create GUI
        self.root = tk.Tk()
        self.root.title("Payroll Mapping Tool")
        self.root.geometry("900x700")
        
        self._create_widgets()
        self._load_existing_mappings()
    
    def _create_widgets(self):
        """Create GUI widgets."""
        # Title
        tk.Label(self.root, text="Define Payroll Mappings",
                font=('Arial', 16, 'bold')).pack(pady=10)
        
        # Scrollable frame for destination columns
        self._create_scrollable_frame()
        
        # Execute button
        tk.Button(self.root, text="Execute Mapping",
                 command=self._execute_mapping,
                 bg='green', fg='white',
                 font=('Arial', 12, 'bold'),
                 width=20, height=2).pack(pady=20)
    
    def _create_scrollable_frame(self):
        """Create scrollable frame with destination columns."""
        frame = tk.Frame(self.root)
        frame.pack(fill=tk.BOTH, expand=True, padx=20, pady=10)
        
        canvas = tk.Canvas(frame)
        scrollbar = tk.Scrollbar(frame, orient="vertical",
                                command=canvas.yview)
        scrollable_frame = tk.Frame(canvas)
        
        scrollable_frame.bind(
            "<Configure>",
            lambda e: canvas.configure(scrollregion=canvas.bbox("all"))
        )
        
        canvas.create_window((0, 0), window=scrollable_frame,
                           anchor="nw")
        canvas.configure(yscrollcommand=scrollbar.set)
        
        # Create row for each destination column
        for i, dest_col in enumerate(CONFIG["target_columns"]):
            self._create_destination_row(scrollable_frame, i, dest_col)
        
        canvas.pack(side="left", fill="both", expand=True)
        scrollbar.pack(side="right", fill="y")
    
    def _create_destination_row(self, parent, row, dest_label):
        """Create row for one destination column."""
        # Label
        tk.Label(parent, text=dest_label,
                font=('Arial', 11), width=25,
                anchor='w').grid(row=row, column=0, padx=5, pady=5)
        
        # Formula entry (read-only)
        entry = tk.Entry(parent, width=40, state='readonly',
                        font=('Arial', 10))
        entry.grid(row=row, column=1, padx=5, pady=5)
        self.formula_entries[dest_label] = entry
        
        # Define button
        btn = tk.Button(parent, text="Define Formula",
                       command=lambda: self._define_formula(dest_label),
                       width=15)
        btn.grid(row=row, column=2, padx=5, pady=5)
    
    def _define_formula(self, dest_label):
        """Open dialog to define formula for destination column."""
        formula = simpledialog.askstring(
            "Define Formula",
            f"Enter formula for '{dest_label}':\n\n"
            f"Examples:\n"
            f"  +SALAIRE_BASE label\n"
            f"  +BASE label +PRIME label\n"
            f"  +BRUT label -COT label\n\n"
            f"Current: {self.mappings.get(dest_label, '(not defined)')}"
        )
        
        if formula:
            self.mappings[dest_label] = formula
            self._update_formula_display(dest_label)
    
    def _execute_mapping(self):
        """Execute all mappings."""
        if not self.mappings:
            messagebox.showwarning("No Mappings",
                                  "Define at least one mapping first")
            return
        
        # Show progress
        progress = tk.Toplevel(self.root)
        progress.title("Mapping in Progress")
        tk.Label(progress, text="Processing...",
                font=('Arial', 12)).pack(padx=50, pady=20)
        progress_bar = ttk.Progressbar(progress, mode='indeterminate')
        progress_bar.pack(padx=20, pady=10)
        progress_bar.start()
        
        # Execute in background thread
        def run_mapping():
            try:
                from src.features.payroll_processing.mapping_engine import (
                    execute_all_mappings
                )
                execute_all_mappings(self.source_wb_path, self.mappings)
                
                # Close progress
                progress_bar.stop()
                progress.destroy()
                
                # Success message
                messagebox.showinfo("Success",
                                   "Mapping completed successfully!")
            except Exception as e:
                progress_bar.stop()
                progress.destroy()
                messagebox.showerror("Error",
                                    f"Mapping failed:\n{str(e)}")
        
        import threading
        threading.Thread(target=run_mapping, daemon=True).start()
    
    def run(self):
        """Start GUI main loop."""
        self.root.mainloop()
\end{lstlisting}

\section{Mapping Engine Modules}

\subsection{src/features/payroll\_processing/mapping\_engine.py}

\textbf{Purpose:} Core mapping engine that applies formulas to transform data.

\textbf{Key Functions:}

\begin{lstlisting}[style=pythonstyle]
@dataclass
class MappingRule:
    """Represents a mapping rule."""
    dest_label_raw: str
    dest_label_norm: str
    mapping_type: str
    keys_ops: List[Tuple[str, str, str]]  # [(op, key, type), ...]
    
    def apply(self, sheet, employee_col, row_lookup):
        """Apply formula to calculate value."""
        result = 0
        
        for operator, key, key_type in self.keys_ops:
            # Find source row
            if key_type == "code":
                row = row_lookup.get_by_code(key)
            else:
                row = row_lookup.get_by_label(key)
            
            if row is None:
                continue
            
            # Get value
            value = sheet.range(row, employee_col).value or 0
            
            # Apply operator
            if operator == '+':
                result += value
            elif operator == '-':
                result -= value
            elif operator == '*':
                result *= value
            elif operator == '/':
                result = result / value if value != 0 else 0
        
        return result


class RowLookup:
    """Fast lookup for source rows."""
    
    def __init__(self, sheet, code_col=0, label_col=2):
        self.code_to_row = {}
        self.label_to_row = {}
        
        # Build dictionaries
        codes = sheet.range("A:A").value
        labels = sheet.range("B:B").value
        
        for i, (code, label) in enumerate(zip(codes, labels), start=1):
            if i < 8:  # Skip header rows
                continue
            
            if code:
                self.code_to_row[str(code).strip()] = i
            
            if label:
                norm = _norm_label(str(label))
                self.label_to_row[norm] = i
    
    def get_by_code(self, code):
        return self.code_to_row.get(str(code).strip())
    
    def get_by_label(self, label):
        return self.label_to_row.get(_norm_label(label))


def execute_all_mappings(source_wb_path, mappings):
    """Execute all mappings."""
    import xlwings as xw
    
    # Open workbooks
    source_wb = xw.Book(source_wb_path)
    target_wb = xw.Book()  # New workbook
    
    # Convert mappings to rules
    rules = []
    for dest_label, formula in mappings.items():
        keys_ops = parse_formula(formula)
        rule = MappingRule(
            dest_label_raw=dest_label,
            dest_label_norm=_norm_label(dest_label),
            mapping_type="formula",
            keys_ops=keys_ops
        )
        rules.append(rule)
    
    # Save rules to CustomMap sheet
    save_mappings(source_wb, rules)
    
    # Process each month sheet
    for month_name in CONFIG["month_sheets"]:
        try:
            src_sheet = source_wb.sheets[month_name]
        except:
            continue
        
        # Create destination sheet
        dest_sheet = target_wb.sheets.add(month_name)
        
        # Apply mappings
        apply_mappings_one_month(src_sheet, dest_sheet, rules)
    
    # Save target workbook
    target_path = CONFIG["paths"]["project_root"] / "data" / "data_sorted.xlsx"
    target_wb.save(target_path)
    target_wb.close()
    source_wb.close()
\end{lstlisting}

\section{Analytics Modules}

\subsection{src/features/analytics/dashboard\_gui.py}

\textbf{Purpose:} Interactive dashboard with 10 visualizations.

\textbf{Size:} 630+ lines

\textbf{Key Methods:} See Chapter 6 (Visualization Dashboard) for detailed documentation of all 10 plots.

\begin{lstlisting}[style=pythonstyle]
class PayrollDashboard:
    """Main dashboard class."""
    
    def __init__(self, data_path):
        self.df = pd.read_csv(data_path)
        self.df['month'] = pd.to_datetime(self.df['month'])
        self.root = tk.Tk()
        self._create_widgets()
    
    # 10 plotting methods
    def plot_total_cost_trend(self): ...
    def plot_cost_forecast(self): ...
    def plot_cost_per_employee(self): ...
    def plot_cost_breakdown(self): ...
    def plot_cost_distribution(self): ...
    def plot_top_employees(self): ...
    def plot_gross_vs_net(self): ...
    def plot_headcount_trend(self): ...
    def plot_burn_rate(self): ...
    def plot_yoy_comparison(self): ...
    
    def run(self):
        self.root.mainloop()
\end{lstlisting}

\subsection{src/features/analytics/regression\_ar1.py to ar4.py}

\textbf{Purpose:} Standalone AR model implementations for testing.

\textbf{Key Functions:}

\begin{lstlisting}[style=pythonstyle]
def run_ar2_model(data_path, forecast_steps=6):
    """Run AR(2) model and save results."""
    df = pd.read_csv(data_path)
    df['month'] = pd.to_datetime(df['month'])
    
    monthly = df.groupby('month')['salaire_brut'].sum()
    
    # Fit AR(2)
    from statsmodels.tsa.ar_model import AutoReg
    model = AutoReg(monthly, lags=2)
    model_fit = model.fit()
    
    # Forecast
    predictions = model_fit.forecast(steps=forecast_steps)
    
    # Save results
    results = pd.DataFrame({
        'month': range(1, forecast_steps+1),
        'prediction': predictions.values
    })
    results.to_csv('outputs/results_ar2_predictions.csv', index=False)
    
    # Employee-level metrics
    emp_metrics = calculate_employee_metrics(df, predictions)
    emp_metrics.to_csv('outputs/results_ar2_emp_metrics.csv', index=False)
\end{lstlisting}

\section{Script Files}

\subsection{scripts/run\_excel.py}

\textbf{Purpose:} Launch mapping GUI.

\begin{lstlisting}[style=pythonstyle]
"""Launch Excel mapping GUI."""

import sys
from pathlib import Path

# Add project root to path
sys.path.insert(0, str(Path(__file__).parent.parent))

from src.excel_tools.import_dialog import select_workbook_dialog
from src.excel_tools.mapping_gui import MappingGUI

def main():
    # Select source workbook
    wb_path = select_workbook_dialog()
    
    if not wb_path:
        print("No file selected. Exiting.")
        return
    
    # Launch GUI
    gui = MappingGUI(wb_path)
    gui.run()

if __name__ == '__main__':
    main()
\end{lstlisting}

\subsection{scripts/run\_dashboard.py}

\textbf{Purpose:} Launch analytics dashboard.

\begin{lstlisting}[style=pythonstyle]
"""Launch analytics dashboard."""

import sys
from pathlib import Path

sys.path.insert(0, str(Path(__file__).parent.parent))

from src.features.analytics.dashboard_gui import PayrollDashboard
from src.core.config import CONFIG

def main():
    data_path = CONFIG["output"]["payroll_long"]
    
    if not Path(data_path).exists():
        print(f"ERROR: Data file not found: {data_path}")
        print("Run data conversion first.")
        return
    
    dashboard = PayrollDashboard(data_path)
    dashboard.run()

if __name__ == '__main__':
    main()
\end{lstlisting}

\subsection{scripts/run\_forecast.py}

\textbf{Purpose:} Run AR(2) forecasting standalone.

\begin{lstlisting}[style=pythonstyle]
"""Run AR(2) forecasting."""

import sys
from pathlib import Path

sys.path.insert(0, str(Path(__file__).parent.parent))

from src.core.config import CONFIG
from src.features.analytics.regression_ar2 import run_ar2_model

def main():
    data_path = CONFIG["output"]["payroll_long"]
    
    print("Running AR(2) forecasting...")
    run_ar2_model(data_path, forecast_steps=6)
    print("Forecast complete. Results saved to outputs/")

if __name__ == '__main__':
    main()
\end{lstlisting}

\subsection{scripts/run\_load\_data.py}

\textbf{Purpose:} Convert Excel to long-format CSV.

\begin{lstlisting}[style=pythonstyle]
"""Convert mapped Excel to long-format CSV."""

import sys
from pathlib import Path

sys.path.insert(0, str(Path(__file__).parent.parent))

from src.core.multi_sheet_loader import load_all_months_to_long
from src.core.config import CONFIG

def main():
    source_path = CONFIG["paths"]["project_root"] / "data" / "data_sorted.xlsx"
    output_path = CONFIG["output"]["payroll_long"]
    
    print(f"Loading data from: {source_path}")
    df = load_all_months_to_long(str(source_path))
    
    print(f"Saving to: {output_path}")
    df.to_csv(output_path, index=False)
    
    print(f"Complete! {len(df)} rows saved.")
    print(f"Columns: {list(df.columns)}")

if __name__ == '__main__':
    main()
\end{lstlisting}

\section{Test Files}

\subsection{tests/test\_monthly\_mapping\_engine.py}

\textbf{Purpose:} Unit tests for mapping engine.

\begin{lstlisting}[style=pythonstyle]
"""Tests for mapping engine."""

import unittest
from src.core.normalization import _norm_label
from src.features.payroll_processing.mapping_engine import (
    MappingRule, RowLookup, parse_formula
)

class TestNormalization(unittest.TestCase):
    def test_norm_label(self):
        self.assertEqual(_norm_label("SALAIRE BASE"), "salaire base")
        self.assertEqual(_norm_label("Salaire  Base"), "salaire base")
        self.assertEqual(_norm_label("Côtisation"), "cotisation")

class TestFormulaParser(unittest.TestCase):
    def test_single_term(self):
        result = parse_formula("+SALAIRE_BASE label")
        self.assertEqual(result, [("+", "SALAIRE_BASE", "label")])
    
    def test_multiple_terms(self):
        result = parse_formula("+BASE label +PRIME label")
        self.assertEqual(len(result), 2)

class TestMappingRule(unittest.TestCase):
    def test_apply_addition(self):
        # Mock objects needed
        pass  # Implementation details

if __name__ == '__main__':
    unittest.main()
\end{lstlisting}

\section{File Responsibilities Summary}

\begin{table}[H]
\centering
\footnotesize
\begin{tabularx}{\textwidth}{lX}
\toprule
\textbf{File} & \textbf{Responsibility} \\
\midrule
config.py & Centralized configuration \\
normalization.py & Label normalization for fuzzy matching \\
data\_loader.py & Load CSV/Excel data \\
multi\_sheet\_loader.py & Combine month sheets to long format \\
import\_dialog.py & File selection dialog \\
mapping\_gui.py & Main mapping GUI (630 lines) \\
mapping\_engine.py & Core mapping logic, formula application \\
dashboard\_gui.py & Interactive visualization dashboard \\
regression\_ar1-ar4.py & AR model implementations \\
regression\_pooled\_fe.py & Panel regression with fixed effects \\
regression\_unpooled.py & Separate regressions per employee \\
visualize\_all\_models.py & Compare all model results \\
plot\_salary\_field.py & Individual field plotting \\
generate\_synthetic\_payroll.py & Synthetic data generation \\
run\_excel.py & Launch mapping GUI (CLI) \\
run\_dashboard.py & Launch dashboard (CLI) \\
run\_forecast.py & Run forecasting standalone \\
run\_load\_data.py & Excel to CSV conversion \\
test\_*.py & Unit tests \\
\bottomrule
\end{tabularx}
\caption{File Responsibilities}
\end{table}

\section{Dependencies}

\subsection{requirements.txt}

\begin{lstlisting}
pandas>=1.5.0
numpy>=1.23.0
matplotlib>=3.6.0
xlwings>=0.28.0
statsmodels>=0.14.0
scikit-learn>=1.1.0
openpyxl>=3.0.0
python-dateutil>=2.8.0
\end{lstlisting}

\subsection{Package Purposes}

\begin{table}[H]
\centering
\begin{tabular}{ll}
\toprule
\textbf{Package} & \textbf{Purpose} \\
\midrule
pandas & Data manipulation, DataFrames \\
numpy & Numerical operations \\
matplotlib & Plotting and visualization \\
xlwings & Excel file I/O \\
statsmodels & AR models, time series analysis \\
scikit-learn & Machine learning utilities \\
openpyxl & Excel file reading (backend) \\
python-dateutil & Date parsing and manipulation \\
\bottomrule
\end{tabular}
\caption{Python Package Dependencies}
\end{table}

\section{VBA Code Reference}

\subsection{Control\_Panel\_New.xlsm Module}

\textbf{Module Name:} Module1

\textbf{Total Buttons:} 5

\textbf{Key Subroutines:}

\begin{lstlisting}[style=vbastyle]
' Button 1: Import & Map Data
Sub ImportAndMapData()
    ' ... (see Chapter 4 for full code)
End Sub

' Button 2: Open Output Folder
Sub OpenOutputFolder()
    Shell "explorer " & Chr(34) & PROJECT_PATH & "outputs" & Chr(34), vbNormalFocus
End Sub

' Button 4: Visualize Data
Sub VisualizeData()
    Shell Chr(34) & PYTHON_PATH & Chr(34) & " " & _
          Chr(34) & PROJECT_PATH & "scripts\run_dashboard.py" & Chr(34), vbNormalFocus
End Sub

' Button 5: Open Mapped Workbook
Sub OpenMappedWorkbook()
    Shell "explorer " & Chr(34) & PROJECT_PATH & "data\data_sorted.xlsx" & Chr(34), vbNormalFocus
End Sub
\end{lstlisting}

\section{Key Algorithms}

\subsection{Formula Parser}

\begin{algorithm}[H]
\caption{Parse Formula String}
\begin{algorithmic}[1]
\Require Formula string (e.g., "+BASE label +PRIME label")
\Ensure List of (operator, key, type) tuples

\State terms $\gets$ split formula by whitespace
\State result $\gets$ empty list
\State i $\gets$ 0

\While{i < len(terms)}
    \State operator $\gets$ terms[i][0]  \Comment{First character: +, -, *, /}
    \State key $\gets$ terms[i][1:]  \Comment{Rest of token}
    \State type $\gets$ terms[i+1]  \Comment{Next token: "code" or "label"}
    
    \State result.append((operator, key, type))
    \State i $\gets$ i + 2
\EndWhile

\State \Return result
\end{algorithmic}
\end{algorithm}

\subsection{Employee Block Detection}

\begin{algorithm}[H]
\caption{Detect Employee Blocks}
\begin{algorithmic}[1]
\Require Source sheet, ID row number
\Ensure List of (employee\_id, column\_number) tuples

\State row\_values $\gets$ read row from sheet
\State blocks $\gets$ empty list

\For{col\_idx, value in enumerate(row\_values)}
    \If{col\_idx $\leq$ 2}
        \State continue  \Comment{Skip code/label columns}
    \EndIf
    
    \If{value is not empty}
        \State emp\_id $\gets$ value.strip()
        \State blocks.append((emp\_id, col\_idx))
    \EndIf
\EndFor

\State \Return blocks
\end{algorithmic}
\end{algorithm}

\subsection{Mapping Application}

See Chapter 5, Section 5.1 for detailed algorithm.
